% ============================================================
%  GUÍA 9: Definición de Derivada y Reglas Básicas
%  Curso de Introducción a la Universidad — Precálculo y Cálculo I
%  Autor: Ing. Gabriel Astudillo
%  Clase cubierta: 9
% ============================================================

\documentclass[12pt, letterpaper]{article}

% ── Paquetes ─────────────────────────────────────────────────
\usepackage{mathwizards-guia}
\mwguia{Guía 9 — Definición de Derivada y Reglas Básicas}{Ing. Gabriel Astudillo $\cdot$ Curso Preuniversitario}

\begin{document}

% ── Portada / Título ─────────────────────────────────────────
\mwportada{Guía de Estudio 9}{Definición de Derivada y Reglas Básicas}{Cociente Incremental, Recta Tangente, Reglas de Potencia, Suma, Producto y Cociente}

\vspace{4pt}
\begin{cuadroinfo}
  \small
  \textbf{Presentación:} Esta guía introduce el concepto de derivada a partir del cociente incremental, la ecuación de la recta tangente y las reglas básicas de derivación: potencia, suma, producto y cociente. La práctica intensiva con ejercicios resueltos y propuestos —con su clave— consolida la técnica algebraica que exige el cálculo diferencial.
  \\[4pt]
  \textbf{Nivel de Dificultad:}\quad
  \facil{} Básico / Directo \quad
  \medio{} Intermedio \quad
  \dificil{} Desafiante \quad
  \muyDificil{} Muy Difícil / Reto
\end{cuadroinfo}

\vspace{8pt}

\section{Definición de Derivada y Recta Tangente}
% ============================================================

\subsection{Conceptos fundamentales}

\begin{cuadroteoria}
  \textbf{Definición de derivada en $x = a$:}
  $$f'(a) = \lim_{h \to 0} \frac{f(a + h) - f(a)}{h}$$
  o equivalentemente:
  $$f'(a) = \lim_{x \to a} \frac{f(x) - f(a)}{x - a}$$

  \textbf{Definición general (función derivada):}
  $$f'(x) = \lim_{h \to 0} \frac{f(x + h) - f(x)}{h}$$

  \textbf{Interpretación geométrica:} $f'(a)$ es la pendiente de la recta tangente a la gráfica de $f$ en el punto $(a, f(a))$.
\end{cuadroteoria}

\newpage

\begin{cuadrosolucion}
  \textbf{Recta tangente en $(a, f(a))$:}
  $$y - f(a) = f'(a)(x - a)$$

  \textbf{Notaciones para la derivada:}
  $$f'(x) \qquad y' \qquad \frac{dy}{dx} \qquad \frac{d}{dx}[f(x)]$$
\end{cuadrosolucion}

\begin{cuadroadvertencia}
  Si el límite no existe, la función \textbf{no es derivable} en ese punto. La continuidad es necesaria pero no suficiente: por ejemplo, $f(x) = |x|$ es continua en $x = 0$ pero no derivable ahí (pico).
\end{cuadroadvertencia}

\subsection{Ejercicios resueltos}

\textbf{Ejemplo R1.} Usando la definición de derivada, calcula $f'(x)$ para $f(x) = x^2$.

\begin{cuadrosolucion}
  \textbf{Solución:}
  \begin{align*}
    f'(x) & = \lim_{h \to 0} \frac{(x + h)^2 - x^2}{h}       \\
          & = \lim_{h \to 0} \frac{x^2 + 2xh + h^2 - x^2}{h} \\
          & = \lim_{h \to 0} \frac{2xh + h^2}{h}             \\
          & = \lim_{h \to 0} (2x + h) = 2x.
  \end{align*}
  Por tanto: $f'(x) = 2x$.
\end{cuadrosolucion}

\textbf{Ejemplo R2.} Usando la definición, calcula la derivada de $f(x) = \dfrac{1}{x}$.

\begin{cuadrosolucion}
  \textbf{Solución:}
  \begin{align*}
    f'(x) & = \lim_{h \to 0} \frac{\frac{1}{x + h} - \frac{1}{x}}{h} \\
          & = \lim_{h \to 0} \frac{\frac{x - (x + h)}{x(x + h)}}{h}  \\
          & = \lim_{h \to 0} \frac{\frac{-h}{x(x + h)}}{h}           \\
          & = \lim_{h \to 0} \frac{-1}{x(x + h)} = -\frac{1}{x^2}.
  \end{align*}
  Por tanto: $f'(x) = -\dfrac{1}{x^2}$.
\end{cuadrosolucion}

\textbf{Ejemplo R3.} Encuentra la ecuación de la recta tangente a $f(x) = x^2$ en $x = 3$.

\begin{cuadrosolucion}
  \textbf{Solución:}
  \begin{enumerate}
    \item Punto: $f(3) = 9$, así que el punto es $(3, 9)$.
    \item Pendiente: $f'(3) = 2(3) = 6$.
    \item Ecuación: $y - 9 = 6(x - 3) \Rightarrow y = 6x - 18 + 9 \Rightarrow y = 6x - 9$.
  \end{enumerate}
\end{cuadrosolucion}

\textbf{Ejemplo R4.} Usando la definición, calcula $f'(x)$ para $f(x) = \sqrt{x}$.

\begin{cuadrosolucion}
  \textbf{Solución:}
  \begin{align*}
    f'(x) & = \lim_{h \to 0} \frac{\sqrt{x + h} - \sqrt{x}}{h} \cdot \frac{\sqrt{x + h} + \sqrt{x}}{\sqrt{x + h} + \sqrt{x}} \\
          & = \lim_{h \to 0} \frac{(x + h) - x}{h(\sqrt{x + h} + \sqrt{x})}                                                  \\
          & = \lim_{h \to 0} \frac{h}{h(\sqrt{x + h} + \sqrt{x})}                                                            \\
          & = \lim_{h \to 0} \frac{1}{\sqrt{x + h} + \sqrt{x}} = \frac{1}{2\sqrt{x}}.
  \end{align*}
  Por tanto: $f'(x) = \dfrac{1}{2\sqrt{x}}$ para $x > 0$.
\end{cuadrosolucion}

\textbf{Ejemplo R5.} Determina si $f(x) = |x|$ es derivable en $x = 0$.

\begin{cuadrosolucion}
  \textbf{Solución:} Calculamos el límite por la izquierda y por la derecha:
  \begin{itemize}
    \item Por la derecha: $\lim_{h \to 0^+} \frac{|0 + h| - |0|}{h} = \lim_{h \to 0^+} \frac{h}{h} = 1$.
    \item Por la izquierda: $\lim_{h \to 0^-} \frac{|h| - 0}{h} = \lim_{h \to 0^-} \frac{-h}{h} = -1$.
  \end{itemize}
  Los límites laterales son distintos, por lo que $f'(0)$ \textbf{no existe}. La función tiene un pico en el origen.
\end{cuadrosolucion}

\newpage

\subsection{Ejercicios propuestos}

\begin{enumerate}[series=ejercicios, leftmargin=*]
  \item \facil{} Usando la definición, calcula $f'(x)$ para $f(x) = 3x^2$.
  \item \facil{} Usando la definición, calcula $f'(x)$ para $f(x) = 2x + 1$.
  \item \facil{} Encuentra la pendiente de la recta tangente a $f(x) = x^2$ en $x = -2$.
  \item \facil{} Encuentra la ecuación de la recta tangente a $f(x) = x^2$ en $x = 1$.
  \item \medio{} Usando la definición, calcula $f'(x)$ para $f(x) = x^3$.
  \item \medio{} Usando la definición, calcula $f'(x)$ para $f(x) = x^2 + 3x$.
  \item \medio{} Usando la definición, calcula $f'(x)$ para $f(x) = \dfrac{1}{x^2}$.
  \item \medio{} Determina si $f(x) = \begin{cases} x^2 & \text{si } x \le 0 \\ 2x & \text{si } x > 0 \end{cases}$ es derivable en $x = 0$.
  \item \dificil{} Usando la definición, calcula $f'(x)$ para $f(x) = \sqrt{x + 2}$.
  \item \dificil{} Usando la definición, calcula $f'(x)$ para $f(x) = \cos x$. (Pista: usa $\cos(a + b) = \cos a \cos b - \sin a \sin b$ y los límites $\lim_{h \to 0} \frac{\sin h}{h} = 1$ y $\lim_{h \to 0} \frac{1 - \cos h}{h} = 0$.)
  \item \dificil{} Encuentra los puntos de la gráfica de $f(x) = x^2 - 4x + 3$ donde la recta tangente es horizontal.
  \item \dificil{} Encuentra la ecuación de la recta tangente a $f(x) = \sqrt{x}$ que sea paralela a la recta $y = \dfrac{x}{4} + 2$.
\end{enumerate}

\newpage

% ============================================================

\section{Reglas Básicas de Derivación}
% ============================================================

\subsection{Reglas fundamentales}

\begin{cuadroteoria}
  \textbf{Regla de la constante:} $\dfrac{d}{dx}[c] = 0$

  \textbf{Regla de la potencia:} $\dfrac{d}{dx}[x^n] = n x^{n-1}$

  \textbf{Regla de la constante por función:} $\dfrac{d}{dx}[c \cdot f(x)] = c \cdot f'(x)$

  \textbf{Regla de la suma/resta:} $\dfrac{d}{dx}[f(x) \pm g(x)] = f'(x) \pm g'(x)$

  \textbf{Regla del producto:} $(fg)' = f'g + fg'$

  \textbf{Regla del cociente:} $\left(\dfrac{f}{g}\right)' = \dfrac{f'g - fg'}{g^2}$

  \textbf{Derivada de la función exponencial:} $\dfrac{d}{dx}[e^x] = e^x$

  \textbf{Derivada del logaritmo natural:} $\dfrac{d}{dx}[\ln x] = \dfrac{1}{x}$ (para $x > 0$)

  \textbf{Derivadas de exponenciales y logaritmos generales:}
  $$\frac{d}{dx}[a^x] = a^x \ln a \qquad\qquad \frac{d}{dx}[\log_a x] = \frac{1}{x \ln a}$$
\end{cuadroteoria}

\begin{cuadrosolucion}
  \textbf{Derivadas trigonométricas básicas:}
  \begin{align*}
    \frac{d}{dx}[\sin x] & = \cos x        & \frac{d}{dx}[\cos x] & = -\sin x        \\
    \frac{d}{dx}[\tan x] & = \sec^2 x      & \frac{d}{dx}[\cot x] & = -\csc^2 x      \\
    \frac{d}{dx}[\sec x] & = \sec x \tan x & \frac{d}{dx}[\csc x] & = -\csc x \cot x
  \end{align*}
\end{cuadrosolucion}

\begin{cuadroadvertencia}
  \textbf{Para no equivocarte:}
  \begin{itemize}[leftmargin=*]
    \item $\dfrac{d}{dx}[x^n]$: baja el exponente y resta 1. Ej: $\dfrac{d}{dx}[x^5] = 5x^4$.
    \item $\dfrac{d}{dx}[e^x]$: ¡es $e^x$! No lo confundas con la regla de la potencia.
    \item La regla del producto NO es $(fg)' = f' \cdot g'$. Debes usar $f'g + fg'$.
    \item La regla del cociente: ``abajo al cuadrado, y arriba: la derivada del primero por el segundo, menos el primero por la derivada del segundo''.
  \end{itemize}
\end{cuadroadvertencia}

\newpage

\subsection{Ejercicios resueltos}

\textbf{Ejemplo R1.} Calcula la derivada de $f(x) = 5x^4 - 3x^2 + 2x - 7$.

\begin{cuadrosolucion}
  \textbf{Solución:} Aplicamos la regla de la potencia término a término:
  $$f'(x) = 5(4)x^3 - 3(2)x + 2 = 20x^3 - 6x + 2.$$
  La constante $-7$ se deriva a $0$.
\end{cuadrosolucion}

\textbf{Ejemplo R2.} Calcula la derivada de $g(x) = e^x + \ln x + \sin x$.

\begin{cuadrosolucion}
  \textbf{Solución:}
  $$g'(x) = e^x + \frac{1}{x} + \cos x.$$
\end{cuadrosolucion}

\textbf{Ejemplo R3.} Calcula la derivada de $h(x) = x^2 e^x$ (regla del producto).

\begin{cuadrosolucion}
  \textbf{Solución:} Identificamos $f = x^2$, $g = e^x$. Entonces $f' = 2x$, $g' = e^x$:
  $$h'(x) = (2x)(e^x) + (x^2)(e^x) = e^x(2x + x^2) = x e^x (2 + x).$$
\end{cuadrosolucion}

\textbf{Ejemplo R4.} Calcula la derivada de $r(x) = \dfrac{x^3}{x^2 + 1}$ (regla del cociente).

\begin{cuadrosolucion}
  \textbf{Solución:} $f = x^3$, $g = x^2 + 1$; $f' = 3x^2$, $g' = 2x$:
  $$r'(x) = \frac{(3x^2)(x^2 + 1) - (x^3)(2x)}{(x^2 + 1)^2} = \frac{3x^4 + 3x^2 - 2x^4}{(x^2 + 1)^2} = \frac{x^4 + 3x^2}{(x^2 + 1)^2} = \frac{x^2(x^2 + 3)}{(x^2 + 1)^2}.$$
\end{cuadrosolucion}

\textbf{Ejemplo R5.} Calcula la derivada de $f(x) = \tan x \cdot \ln x$.

\begin{cuadrosolucion}
  \textbf{Solución:} Por producto con $f = \tan x$, $g = \ln x$:
  $$f'(x) = \sec^2 x \cdot \ln x + \tan x \cdot \frac{1}{x}.$$
\end{cuadrosolucion}

\textbf{Ejemplo R6.} Calcula la derivada de $y = \dfrac{\sin x}{x}$.

\begin{cuadrosolucion}
  \textbf{Solución:} Por cociente con $f = \sin x$, $g = x$:
  $$y' = \frac{\cos x \cdot x - \sin x \cdot 1}{x^2} = \frac{x \cos x - \sin x}{x^2}.$$
\end{cuadrosolucion}

\newpage

\subsection{Ejercicios propuestos}

\begin{enumerate}[resume=ejercicios, leftmargin=*]
  \item \facil{} Calcula: $f'(x)$ si $f(x) = 3x^5 - 4x^3 + x - 9$.
  \item \facil{} Calcula: $g'(x)$ si $g(x) = \dfrac{1}{x^3}$ (pista: reescribe como $x^{-3}$).
  \item \facil{} Calcula: $h'(x)$ si $h(x) = 7e^x + 2\ln x$.
  \item \facil{} Calcula: $y'$ si $y = 2\sin x - \cos x$.
  \item \medio{} Calcula: $f'(x)$ si $f(x) = x^3 \sin x$.
  \item \medio{} Calcula: $g'(x)$ si $g(x) = \dfrac{x^2}{x - 1}$.
  \item \medio{} Calcula: $h'(x)$ si $h(x) = (x^2 + 1) \ln x$.
  \item \medio{} Calcula: $y'$ si $y = \dfrac{e^x}{x}$.
  \item \dificil{} Calcula: $f'(x)$ si $f(x) = \dfrac{x^2 + 1}{x^2 - 1}$.
  \item \dificil{} Calcula: $g'(x)$ si $g(x) = x^2 e^x \sin x$ (aplica producto dos veces).
  \item \dificil{} Calcula: $h'(x)$ si $h(x) = \dfrac{\tan x}{e^x}$.
  \item \dificil{} Encuentra todos los puntos donde la recta tangente a $f(x) = x^3 - 3x$ es horizontal.
\end{enumerate}

% ============================================================

% ──────────────────────────────────────────────────────────
%  NUEVOS EJERCICIOS PROPUESTOS (26) — INSERCIÓN MANUAL
% ──────────────────────────────────────────────────────────
\subsection{Ejercicios propuestos: Recta Tangente y Definición}

\begin{enumerate}[resume=ejercicios, leftmargin=*]
  \item \facil{} Usa la definición de derivada (cociente incremental) para hallar $f'(x)$ si $f(x) = 3x^2$.

  \item \facil{} Deriva usando las reglas básicas: $f(x) = 5x^3 - 2x$.

  \item \medio{} Halla la pendiente de la recta tangente a $y = x^2$ en $x = 3$.

  \item \medio{} Halla la ecuación de la recta tangente a $f(x) = x^2$ en el punto de abscisa $x = 2$.

  \item \medio{} Deriva usando la regla del producto: $f(x) = (2x - 1)(x + 3)$.

  \item \dificil{} Deriva usando la regla del cociente: $f(x) = \dfrac{2x + 1}{x - 3}$.

  \item \dificil{} Halla la ecuación de la recta tangente a $y = \dfrac{1}{x}$ en $x = 2$.

  \item \dificil{} Determina los puntos donde la recta tangente a $y = x^2 - 4x + 3$ es horizontal.

  \item \muyDificil{} Usa la definición de derivada para hallar $f'(x)$ si $f(x) = \sqrt{x}$.
\end{enumerate}

\subsection{Ejercicios propuestos: Reglas Básicas Avanzadas}

\begin{enumerate}[resume=ejercicios, leftmargin=*]
  \item \facil{} Deriva: $f(x) = x^5 - 3x^3 + 2x$.

  \item \facil{} Deriva la función constante: $f(x) = 6$.

  \item \medio{} Deriva usando la regla del producto: $f(x) = (x^2 + 1)(x^3 - 2)$.

  \item \medio{} Deriva usando la regla del cociente: $f(x) = \dfrac{x^2}{x + 1}$.

  \item \medio{} Halla la ecuación de la recta tangente a $y = x^3$ en $x = 1$.

  \item \dificil{} Deriva expandiendo y aplicando el producto: $f(x) = (3x + 2)^2 (x - 1)$.

  \item \dificil{} Halla la ecuación de la recta normal a $y = x^2$ en $x = 1$.

  \item \dificil{} Determina los puntos de $y = \dfrac{x}{x + 1}$ cuya recta tangente tiene pendiente $\dfrac{1}{4}$.

  \item \muyDificil{} Encuentra los puntos de $y = x^3$ donde la recta tangente es paralela a la recta $y = 3x$.
\end{enumerate}

\subsection{Ejercicios propuestos: Retos Integradores}

\begin{enumerate}[resume=ejercicios, leftmargin=*]
  \item \facil{} Deriva: $f(x) = x^7 - 2x^4 + 3x - 1$.

  \item \medio{} Deriva simplificando primero: $f(x) = (x^3 + 1)(x^3 - 1)$.

  \item \medio{} Deriva usando la regla del cociente: $f(x) = \dfrac{x^2 - 1}{x^2 + 1}$.

  \item \medio{} \textbf{Cinemática:} La posición de un móvil es $s(t) = t^2 + 3t$ (en metros). Determina su velocidad en $t = 2$ s.

  \item \dificil{} Calcula $f'(2)$ para $f(x) = \dfrac{x^2 + 1}{x}$.

  \item \dificil{} Halla los puntos de $y = x^2 + 1$ cuya recta tangente pasa por el origen.

  \item \muyDificil{} Usa la definición de derivada para hallar $f'(x)$ si $f(x) = \dfrac{1}{x}$.

  \item \muyDificil{} Determina $a$ y $b$ para que la recta tangente a $y = ax^2 + bx$ en $x = 1$ tenga pendiente 3 y pase por el punto $(1, 5)$.
\end{enumerate}
\newpage
% ============================================================
\ifmwclaves
\section{Clave de Respuestas}
% ============================================================

\subsection*{Sección 1: Definición (Ejercicios 1--12)}

\begin{enumerate}[series=respuestas, leftmargin=*, label=\textbf{\arabic*.}]
  \item $f'(x) = 6x$
  \item $f'(x) = 2$
  \item $f'(-2) = 2(-2) = -4$
  \item Pendiente: $f'(1) = 2$; punto $(1,1)$; recta: $y - 1 = 2(x - 1) \Rightarrow y = 2x - 1$.
  \item Usando $(x+h)^3 = x^3 + 3x^2h + 3xh^2 + h^3$: $f'(x) = 3x^2$.
  \item $f(x+h) - f(x) = (x+h)^2 + 3(x+h) - x^2 - 3x = 2xh + h^2 + 3h$. Dividiendo entre $h$ y tomando límite: $f'(x) = 2x + 3$.
  \item $\dfrac{1}{(x+h)^2} - \dfrac{1}{x^2} = \dfrac{x^2 - (x+h)^2}{x^2(x+h)^2} = \dfrac{-2xh - h^2}{x^2(x+h)^2}$. Dividiendo entre $h$: $\dfrac{-2x - h}{x^2(x+h)^2} \to \dfrac{-2}{x^3}$. Por tanto $f'(x) = -\dfrac{2}{x^3}$.
  \item Por izquierda: $\lim_{h\to 0^-} \frac{h^2}{h} = 0$. Por derecha: $\lim_{h\to 0^+} \frac{2h}{h} = 2$. Son distintos: no es derivable en $x=0$.
  \item Racionalizando: $f'(x) = \dfrac{1}{2\sqrt{x+2}}$.
  \item $\cos(x+h) = \cos x \cos h - \sin x \sin h$. Entonces $\frac{\cos(x+h) - \cos x}{h} = \cos x \frac{\cos h - 1}{h} - \sin x \frac{\sin h}{h} \to \cos x(0) - \sin x(1) = -\sin x$.
  \item Tangente horizontal cuando $f'(x) = 0$: $2x - 4 = 0 \Rightarrow x = 2$, $y = 4 - 8 + 3 = -1$. Punto: $(2, -1)$.
  \item $f'(x) = \frac{1}{2\sqrt{x}}$. La pendiente deseada es $\frac{1}{4}$. $\frac{1}{2\sqrt{x}} = \frac{1}{4} \Rightarrow 2\sqrt{x} = 4 \Rightarrow \sqrt{x} = 2 \Rightarrow x = 4$, $y = 2$. Recta: $y - 2 = \frac{1}{4}(x - 4) \Rightarrow y = \frac{x}{4} + 1$.
\end{enumerate}


\subsection*{Sección 2: Reglas Básicas (Ejercicios 13--24)}

\begin{enumerate}[resume=respuestas, leftmargin=*, label=\textbf{\arabic*.}]
  \item $f'(x) = 15x^4 - 12x^2 + 1$
  \item $g(x) = x^{-3} \Rightarrow g'(x) = -3x^{-4} = -\dfrac{3}{x^4}$
  \item $h'(x) = 7e^x + \dfrac{2}{x}$
  \item $y' = 2\cos x + \sin x$
  \item $f'(x) = 3x^2 \sin x + x^3 \cos x$
  \item $g'(x) = \dfrac{2x(x - 1) - x^2(1)}{(x - 1)^2} = \dfrac{2x^2 - 2x - x^2}{(x - 1)^2} = \dfrac{x^2 - 2x}{(x - 1)^2} = \dfrac{x(x - 2)}{(x - 1)^2}$
  \item $h'(x) = 2x \ln x + (x^2 + 1) \cdot \dfrac{1}{x} = 2x\ln x + x + \dfrac{1}{x}$
  \item $y' = \dfrac{e^x x - e^x}{x^2} = \dfrac{e^x(x - 1)}{x^2}$
  \item $f'(x) = \dfrac{2x(x^2 - 1) - (x^2 + 1)2x}{(x^2 - 1)^2} = \dfrac{2x^3 - 2x - 2x^3 - 2x}{(x^2 - 1)^2} = \dfrac{-4x}{(x^2 - 1)^2}$
  \item Derivamos por producto (tres factores). Podemos agrupar: $g = (x^2 e^x) \sin x = (2x e^x + x^2 e^x)\sin x + x^2 e^x \cos x = x e^x(2 + x)\sin x + x^2 e^x \cos x$.
  \item Por cociente: $h' = \dfrac{\sec^2 x \cdot e^x - \tan x \cdot e^x}{(e^x)^2} = \dfrac{\sec^2 x - \tan x}{e^x}$.
  \item $f'(x) = 3x^2 - 3 = 3(x^2 - 1) = 3(x - 1)(x + 1)$. Ceros: $x = \pm 1$. Puntos: $(1, -2)$ y $(-1, 2)$.
\end{enumerate}


% ──────────────────────────────────────────────────────────
%  NUEVAS RESPUESTAS (26) — INSERCIÓN MANUAL
% ──────────────────────────────────────────────────────────
\subsection*{Sección 3: Recta Tangente y Definición (Ejercicios 25--33)}
\begin{enumerate}[resume=respuestas, leftmargin=*, label=\textbf{\arabic*.}]
  \item $f'(x) = \displaystyle\lim_{h \to 0} \frac{3(x + h)^2 - 3x^2}{h} = 6x$
  \item $f'(x) = 15x^2 - 2$
  \item $m = 2(3) = 6$
  \item $y = 4x - 4$ \quad ($m = 4$, punto $(2, 4)$)
  \item $f'(x) = 4x + 5$
  \item $f'(x) = \dfrac{2(x - 3) - (2x + 1)}{(x - 3)^2} = \dfrac{-7}{(x - 3)^2}$
  \item $y = -\dfrac{x}{4} + 1$ \quad (punto $\left(2, \frac{1}{2}\right)$, pendiente $-\frac{1}{4}$)
  \item $y' = 2x - 4 = 0 \Rightarrow x = 2$; punto $(2, -1)$
  \item $f'(x) = \dfrac{1}{2\sqrt{x}}$
\end{enumerate}

\subsection*{Sección 4: Reglas Básicas Avanzadas (Ejercicios 34--42)}
\begin{enumerate}[resume=respuestas, leftmargin=*, label=\textbf{\arabic*.}]
  \item $f'(x) = 5x^4 - 9x^2 + 2$
  \item $f'(x) = 0$
  \item $f'(x) = 5x^4 + 3x^2 - 4x$
  \item $f'(x) = \dfrac{x^2 + 2x}{(x + 1)^2}$
  \item $y = 3x - 2$
  \item $f'(x) = 27x^2 + 6x - 8$
  \item $y - 1 = -\dfrac{1}{2}(x - 1)$ \quad (pendiente de la normal $-\frac{1}{2}$)
  \item $f'(x) = \dfrac{1}{(x + 1)^2}$; \quad $\dfrac{1}{(x+1)^2} = \dfrac{1}{4} \Rightarrow x = 1$ o $x = -3$; puntos $\left(1, \frac{1}{2}\right)$ y $\left(-3, \frac{3}{2}\right)$
  \item $3x^2 = 3 \Rightarrow x = \pm 1$; puntos $(1, 1)$ y $(-1, -1)$
\end{enumerate}

\subsection*{Sección 5: Retos Integradores (Ejercicios 43--50)}
\begin{enumerate}[resume=respuestas, leftmargin=*, label=\textbf{\arabic*.}]
  \item $f'(x) = 7x^6 - 8x^3 + 3$
  \item $f(x) = x^6 - 1 \Rightarrow f'(x) = 6x^5$
  \item $f'(x) = \dfrac{4x}{(x^2 + 1)^2}$
  \item $v(t) = 2t + 3$; \quad $v(2) = 7$ m/s
  \item $f(x) = x + \dfrac{1}{x}$; \quad $f'(2) = 1 - \dfrac{1}{4} = \dfrac{3}{4}$
  \item Tangente en $(a, a^2 + 1)$: pendiente $2a = \dfrac{a^2 + 1}{a} \Rightarrow a = \pm 1$; puntos $(1, 2)$ y $(-1, 2)$
  \item $f'(x) = \displaystyle\lim_{h \to 0} \frac{\frac{1}{x+h} - \frac{1}{x}}{h} = -\frac{1}{x^2}$
  \item $a + b = 5$ y $2a + b = 3$ $\Rightarrow$ $a = -2$, $b = 7$
\end{enumerate}
\fi
\end{document}
